\documentclass{article}
\usepackage{graphicx} % Required for inserting images
\usepackage{amsfonts}
\usepackage{amsmath}

\title{Apostol Chapter 1 - Exercise 2}
\author{Afonso Vilhena Francisco}
\date{August 2026}

\begin{document}

\maketitle

\textbf{If } $\mathbf{(a,b)=(a,c)=1,}$ \textbf{then } $\mathbf{(a,bc)=1}$

\vspace{0.5cm}

Obviously, $1|a$ and $1|bc$.

Let $e \in \mathbb{Z}$ such that $e|a$ and $e|bc$.

Since $(a,b)=1$, we can write 
$$1=ax+by \;\text{for some integers }x,y,$$
multiplying both sides by $c$ we get
$$c=acx+bcy.$$

From $e|a$ and $e|bc$ we have that there exists $q_1,q_2 \in \mathbb{Z}$ such that
$$a=eq_1 \text{ and } bc=eq_2.$$

Substituing in the previous equation we get

$$c=eq_1cx+eq_2y=e(q_1cx+q_2y),$$

which means that $e|c$. Since $e|a$ and $e|c$ we get that $e|(a,c)$, but $(a,c)=1$, so $e|1$. Therefore $e=\pm1$. In particular, $e|1$. Therefore $(a,bc)=1$

\vspace{1cm}


\end{document}
